\documentclass{article}
\usepackage{amsmath, amssymb}
\usepackage{booktabs}
\usepackage{graphicx}   
\usepackage{xeCJK}
\begin{document}

\title{立定跳远中姿态对成绩影响的力学分析}
\author{孟祥祯}
\date{\today}
\maketitle

\begin{abstract}
立定跳远是一项依赖全身协调发力的运动，姿态调整对跳远成绩有显著提升作用。本文首先建立了一个二刚体（躯干+双腿）模型，描述了人体通过收腿调整姿态、以增加落地前伸量的过程。模型预测的落地前伸量与 Ashby \& Heegaard (2002) 的限制摆臂实验数据吻合。在此基础上，将模型推广至三刚体（躯干+双臂+双腿），进一步分析了手臂运动在腾空阶段对角动量分配的影响。
\end{abstract}

\section{引言}

立定跳远是一项衡量人体爆发力与协调性的基础运动项目。与垂直纵跳不同，立定跳远要求身体在水平方向获得尽可能大的位移，同时落地时保持平衡。实验研究表明，自由摆臂可以显著提升立定跳远成绩。Ashby \& Heegaard (2002) 的实验发现，受试者在自由摆臂时平均跳远距离约为 2.09 m，而手臂受限时仅约为 1.72 m，提升了 21.2\%。其中约 71\% 的增益来自起跳速度的提高，其余来自起跳前质心水平位移的增加和落地时脚超前质心距离的增大。后续的仿真研究（Ashby \& Delp, 2006）进一步指出，手臂摆动通过“前摆储力—后摆控姿”的机制，既增加了系统总能量，又改善了腾空后的角动量分配。

本文首先建立一个二刚体模型（躯干+双腿），该模型能够在将手臂近似与躯干视为一体的前提下，定量描述腾空阶段收腿动作对落地前伸量的贡献。然后，将二刚体模型的预测与 Ashby \& Heegaard (2002) 的实验数据进行对比。最后，将模型推广至三刚体（躯干+双臂+双腿），给出手臂作为独立自由度的定量描述。

\section{立定跳远模型构建}

\subsection{模型定义}
将人体简化为两个刚体：
\begin{itemize}
\item 刚体1（躯干+头+臂）：质量 $m_1$，绕自身质心 $G_1$ 的转动惯量 $I_1$。质心 $G_1$ 到髋关节 $O$ 的距离为 $d_1$，方向沿躯干向上，与水平方向夹角 $\theta_1$。
\item 刚体2（双腿+足）：质量 $m_2$，绕自身质心 $G_2$ 的转动惯量 $I_2$。质心 $G_2$ 到髋关节 $O$ 的距离为 $d_2$，方向沿腿指向脚底，单位向量 $\hat{\mathbf{u}}_2$（见下），与水平方向夹角 $\theta_2$。
\end{itemize}
总质量 $M = m_1 + m_2$，约化质量 $\mu = \dfrac{m_1 m_2}{M}$。

\subsection{运动学关系}
取髋关节 $O$ 为参考点，质心位置：
\begin{align}
\vec{r}_{G_1} &= \vec{r}_O + d_1 \hat{u}_1, \quad \hat{u}_1 = (\cos\theta_1,\sin\theta_1), \\
\vec{r}_{G_2} &= \vec{r}_O + d_2 \hat{u}_2, \quad \hat{u}_2 = (\cos\theta_2,\sin\theta_2), \\
\vec{r}_{12} &= \vec{r}_{G_1} - \vec{r}_{G_2} = d_1\hat{u}_1 - d_2\hat{u}_2.
\end{align}
速度：
\begin{align}
\dot{\vec{r}}_{G_1} &= \dot{\vec{r}}_O + d_1 \dot\theta_1 \hat{k} \times \hat{u}_1, \\
\dot{\vec{r}}_{G_2} &= \dot{\vec{r}}_O + d_2 \dot\theta_2 \hat{k} \times \hat{u}_2, \\
\dot{\vec{r}}_{12} &= d_1 \dot\theta_1 (\hat{k}\times\hat{u}_1) - d_2 \dot\theta_2 (\hat{k}\times\hat{u}_2).
\end{align}
其中 $\hat{k}$ 为垂直于运动平面的单位向量。

计算叉积的 $z$ 分量（标量）：
\begin{align}
(\vec{r}_{12} \times \dot{\vec{r}}_{12})_z = \;& d_1^2 \dot\theta_1 + d_2^2 \dot\theta_2 \nonumber \\
&- d_1 d_2 \cos(\theta_1-\theta_2) (\dot\theta_1 + \dot\theta_2). \label{eq:orbit_exact}
\end{align}
推导利用了恒等式 $(\hat{u}_i \times (\hat{k}\times\hat{u}_j))_z = \hat{u}_i\cdot\hat{u}_j = \cos(\theta_i-\theta_j)$。

\subsection{角动量守恒}
需要说明的是，$L_0=0$ 是一种近似假设。起跳时地面反作用力（GRF）的作用线若通过身体总质心，则角冲量为零。Ashby \& Heegaard (2002) 的实测数据显示，起跳时刻质心与压力中心的水平偏差通常在 5--8\,cm 以内，起跳角度在 $35^\circ$--$40^\circ$ 范围内时 GRF 方向接近质心--足底连线。若起跳技术不理想或姿态偏离较大，$L_0$ 可取非零常值，但推导较为繁琐。本文二刚体模型中，认为初角动量近似为0。

腾空后仅重力作用，对总质心 $G$ 的角动量守恒。设起跳时总角动量为零，则：

\[
L_G = I_1\dot\theta_1 + I_2\dot\theta_2 + \mu\,(\vec{r}_{12} \times \dot{\vec{r}}_{12})_z = 0.
\]
代入式\eqref{eq:orbit_exact}得：
\[
\boxed{I_1\dot\theta_1 + I_2\dot\theta_2 + \mu\left[ d_1^2 \dot\theta_1 + d_2^2 \dot\theta_2 - d_1 d_2 \cos(\theta_1-\theta_2) (\dot\theta_1 + \dot\theta_2) \right] = 0}. \label{eq:angular_momentum_exact}
\]

\subsection{用相对角速度表示}
定义相对角 $\phi = \theta_2 - \theta_1$，则 $\dot\theta_2 = \dot\theta_1 + \dot\phi$。代入式\eqref{eq:angular_momentum_exact}整理得：
\[
\bigl[ I_1 + \mu d_1^2 - \mu d_1 d_2 \cos\phi \bigr] \dot\theta_1
+ \bigl[ I_2 + \mu d_2^2 - \mu d_1 d_2 \cos\phi \bigr] (\dot\theta_1 + \dot\phi) = 0.
\]
解得：
\begin{align}
\dot\theta_1 &= -\,\frac{I_2 + \mu d_2^2 - \mu d_1 d_2 \cos\phi}{I_1+I_2 + \mu(d_1^2+d_2^2) - 2\mu d_1 d_2 \cos\phi}\; \dot\phi, \label{eq:dt1} \\
\dot\theta_2 &= \dot\theta_1 + \dot\phi = \frac{I_1 + \mu d_1^2 - \mu d_1 d_2 \cos\phi}{I_1+I_2 + \mu(d_1^2+d_2^2) - 2\mu d_1 d_2 \cos\phi}\; \dot\phi. \label{eq:dt2}
\end{align}
其中 $\dot\phi(t)$ 由运动员主动控制（例如匀速收腿）。

\subsection{落地前伸量}
脚的位置：
\[
x_B = x_O + l_2 \cos\theta_2.
\]
总质心 $G$ 的水平坐标：
\[
x_G = x_O + \frac{m_1 d_1 \cos\theta_1 + m_2 d_2 \cos\theta_2}{M}.
\]
脚相对于质心的水平偏移为：
\[
\boxed{\Delta x = l_2\cos\theta_2 - \frac{m_1 d_1 \cos\theta_1 + m_2 d_2 \cos\theta_2}{M}}. \label{eq:Dx}
\]
利用式\eqref{eq:dt1}和\eqref{eq:dt2}积分得到 $\theta_1(t),\theta_2(t)$，代入\eqref{eq:Dx}即可得 $\Delta x(t)$。落地时刻 $t=T$ 的 $\Delta x(T)$ 即为通过姿态调整获得的额外跳跃距离。

\subsection{数值求解}
上述方程可采用数值积分求解。人体参数取值如表1所示，腾空时间 $T=0.5\,\text{s}$，收腿角速度 $\dot\phi = 180^\circ/\text{s}$。我编写了Python程序进行数值积分，在 $t=T$ 时的数值由式\eqref{eq:Dx} 给出（结果见第~3.2~节）。

\begin{table}[h]
\centering
\caption{二刚体模型参数}
\label{tab:two_body_params}
\begin{tabular}{lcc}
\toprule
参数 & 符号 & 参考值 \\
\midrule
躯干+头+臂质量 & $m_1$ & 44 kg \\
双腿+足质量 & $m_2$ & 28 kg \\
躯干质心到髋距离 & $d_1$ & 0.35 m \\
双腿质心到髋距离 & $d_2$ & 0.35 m \\
腿长（髋到脚底） & $l_2$ & 0.95 m \\
躯干绕质心转动惯量 & $I_1$ & 1.8 kg\,m$^2$ \\
双腿绕质心转动惯量 & $I_2$ & 2.1 kg\,m$^2$ \\
起跳时躯干倾角 & $\theta_1(0)$ & $60^\circ$ \\
起跳时腿倾角 & $\theta_2(0)$ & $-95^\circ$ \\
收腿角速度 & $\dot\phi$ & $180^\circ/\text{s}$ \\
腾空时间 & $T$ & 0.5 s \\
\bottomrule
\end{tabular}
\end{table}

\section{二刚体模型与实验数据的对比}

上一节建立的二刚体模型仅描述了腾空阶段，起跳参数（起跳速度 $v_0$、起跳角度 $\theta$、起跳时质心高度 $y_0$ 等）需作为外部输入给定。本节将模型的预测与 Ashby \& Heegaard (2002) 的实测数据进行对比。

\subsection{实验数据}

Ashby \& Heegaard (2002) 对三名男性受试者（身高 $1.81\pm0.03$\,m，体重 $72.3\pm13.0$\,kg）在自由摆臂（JFA）和限制摆臂（JRA）两种条件下分别测量了立定跳远的运动学参数。关键实验数据汇总于表~\ref{tab:exp_data}。

\begin{table}[h]
\centering
\caption{自由臂（JFA）与限制臂（JRA）实验数据对比。$r_x$ 为起跳阶段质心的水平位移量；$\Delta x$ 为落地瞬间脚超前质心的水平距离。数据源自 Ashby \& Heegaard (2002)，起跳高度 $y_0$ 根据文中人体构型图及 Ashby \& Delp (2006) 仿真数据估算。}
\label{tab:exp_data}
\begin{tabular}{lccc}
\toprule
参数 & JFA & JRA & 差值 \\
\midrule
跳跃距离 $L$ (m) & $2.09\pm0.03$ & $1.72\pm0.03$ & $+0.37$ \\
起跳速度 $v_0$ (m/s) & $3.32\pm0.03$ & $2.95\pm0.03$ & $+0.37$ \\
起跳角度 $\theta$ ($^\circ$) & $38.6\pm1.1$ & $40.2\pm1.1$ & $-1.6$ \\
起跳质心位移 $r_x$ (m) & $0.57\pm0.01$ & $0.49\pm0.01$ & $+0.08$ \\
起跳质心高度 $y_0$ (m) & $\approx0.93$ & $\approx0.89$ & $\approx+0.04$ \\
落地超前量 $\Delta x$ (m) & $0.31\pm0.01$ & $0.29\pm0.01$ & $+0.02$ \\
\bottomrule
\end{tabular}
\end{table}

\subsection{二刚体模型对落地前伸量的预测}

模型的输出为落地的脚超前质心距离 $\Delta x(t)$，其由腾空阶段的角动量守恒和收腿运动决定。取表~\ref{tab:two_body_params} 中的参数（对应 JRA）及收腿角速度 $\dot\varphi = 180^\circ/\text{s}$，腾空时间 $T=0.5$\,s，数值积分得到 $\Delta x(t)$。

起跳瞬间（$t=0$）：$\theta_1=60^\circ$，$\theta_2=-95^\circ$，相对角 $\phi(0)=-155^\circ$，
\begin{equation}
\Delta x(0)=l_2\cos\theta_2(0)-\frac{m_1 d_1\cos\theta_1(0)+m_2 d_2\cos\theta_2(0)}{M}\approx -0.178\ \text{m}.
\end{equation}
落地时刻（$t=0.5$\,s）：
\begin{equation}
\boxed{\Delta x(T) \approx 0.298\ \text{m}}.
\end{equation}

\subsection{二刚体模型分析与改进}

实验 JRA 的落地前伸量为 $\Delta x=0.29\pm0.01$\,m，二刚体模型预测为 $0.298$\,m，偏差约 $+0.008$\,m（$\ll 3\%$）。二刚体模型在采用杆模型转动惯量估计后，与实验数据基本吻合。

模型局限：
\begin{itemize}
\item 起跳参数（$v_0$、$\theta$、$y_0$ 等）必须由实验提供。
\item 手臂的作用无法推导，需要通过三刚体推广（第~\ref{sec:three_body} 节）来加以解释。
\item 起跳角动量初值 $L_0=0$ 的假设是一种近似，更严谨的处理需要考虑地面反作用力的角冲量。
\end{itemize}

\section{三刚体模型}
\label{sec:three_body}

二刚体模型将双臂与躯干视为单一刚体，仅能描述JRA情形。本节将模型推广为三个刚体：躯干、双臂和双腿。

\subsection{模型定义}

将人体分为三个刚体：
\begin{itemize}
\item \textbf{刚体T（躯干+头）}：质量 $m_T$，绕自身质心的转动惯量 $I_T$。质心 $G_T$ 位于髋关节 $O$ 上方距离 $d_T$ 处。肩关节 $S$ 位于髋关节上方距离 $L_T$ 处（即躯干长度）。
\item \textbf{刚体A（双臂）}：质量 $m_A$，绕自身质心的转动惯量 $I_A$。质心 $G_A$ 位于肩关节 $S$ 沿手臂方向距离 $d_A$ 处。
\item \textbf{刚体L（双腿+足）}：质量 $m_L$，绕自身质心的转动惯量 $I_L$。质心 $G_L$ 位于髋关节 $O$ 沿腿方向距离 $d_L$ 处，腿长（髋至脚底）为 $l_L$。
\end{itemize}

总质量 $M = m_T + m_A + m_L$。各刚体长轴与水平方向的夹角分别为 $\theta_T$、$\theta_A$、$\theta_L$，对应的单位方向向量：
\[
\hat{\mathbf{u}}_T = (\cos\theta_T,\;\sin\theta_T),\quad
\hat{\mathbf{u}}_A = (\cos\theta_A,\;\sin\theta_A),\quad
\hat{\mathbf{u}}_L = (\cos\theta_L,\;\sin\theta_L).
\]
角速度记作 $\omega_T \equiv \dot\theta_T$、$\omega_A \equiv \dot\theta_A$、$\omega_L \equiv \dot\theta_L$。

\begin{table}[h]
\centering
\caption{三刚体模型参数。质量按 Ashby 受试者体型分配（$m_T+m_A = m_1$，$m_L = m_2$），转动惯量由杆模型估算}
\begin{tabular}{lcc}
\toprule
参数 & 符号 & 典型值 \\
\midrule
躯干+头质量           & $m_T$ & 37 kg \\
双臂质量              & $m_A$ & 7 kg \\
双腿+足质量           & $m_L$ & 28 kg \\
躯干质心至髋距离       & $d_T$ & 0.32 m \\
臂质心至肩距离         & $d_A$ & 0.28 m \\
腿质心至髋距离         & $d_L$ & 0.35 m \\
髋至肩距离（躯干长）    & $L_T$ & 0.60 m \\
腿长（髋至脚底）        & $l_L$ & 0.95 m \\
躯干绕质心转动惯量     & $I_T$ & 1.5 kg\,m$^2$ \\
双臂绕质心转动惯量     & $I_A$ & 0.3 kg\,m$^2$ \\
双腿绕质心转动惯量     & $I_L$ & 2.1 kg\,m$^2$ \\
\bottomrule
\end{tabular}
\end{table}

\subsection{运动学关系}

取髋关节 $O$ 为参考点，各刚体质心和肩关节的位矢：
\begin{align}
\mathbf{r}_T &= \mathbf{r}_O + d_T \hat{\mathbf{u}}_T, \label{eq:rT} \\
\mathbf{r}_S &= \mathbf{r}_O + L_T \hat{\mathbf{u}}_T, \label{eq:rS} \\
\mathbf{r}_A &= \mathbf{r}_S + d_A \hat{\mathbf{u}}_A
              = \mathbf{r}_O + L_T \hat{\mathbf{u}}_T + d_A \hat{\mathbf{u}}_A, \label{eq:rA} \\
\mathbf{r}_L &= \mathbf{r}_O + d_L \hat{\mathbf{u}}_L. \label{eq:rL}
\end{align}

总质心位置：
\begin{equation}
\mathbf{r}_G = \mathbf{r}_O + \beta_T \hat{\mathbf{u}}_T + \beta_A \hat{\mathbf{u}}_A + \beta_L \hat{\mathbf{u}}_L, \label{eq:rG_three}
\end{equation}
为简洁，定义系数：
\begin{equation}
\boxed{\beta_T \equiv \frac{m_T d_T + m_A L_T}{M},\quad
       \beta_A \equiv \frac{m_A d_A}{M},\quad
       \beta_L \equiv \frac{m_L d_L}{M}}. \label{eq:beta_def}
\end{equation}

速度关系：
\begin{align}
\mathbf{v}_T &= \mathbf{v}_O + d_T \omega_T (\hat{\mathbf{k}}\times\hat{\mathbf{u}}_T), \label{eq:vT} \\
\mathbf{v}_A &= \mathbf{v}_O + L_T \omega_T (\hat{\mathbf{k}}\times\hat{\mathbf{u}}_T) + d_A \omega_A (\hat{\mathbf{k}}\times\hat{\mathbf{u}}_A), \label{eq:vA} \\
\mathbf{v}_L &= \mathbf{v}_O + d_L \omega_L (\hat{\mathbf{k}}\times\hat{\mathbf{u}}_L), \label{eq:vL} \\
\mathbf{v}_G &= \mathbf{v}_O + \beta_T \omega_T (\hat{\mathbf{k}}\times\hat{\mathbf{u}}_T) + \beta_A \omega_A (\hat{\mathbf{k}}\times\hat{\mathbf{u}}_A) + \beta_L \omega_L (\hat{\mathbf{k}}\times\hat{\mathbf{u}}_L). \label{eq:vG_three}
\end{align}
其中 $\hat{\mathbf{k}}$ 为垂直于运动平面的单位向量（$z$ 轴正方向）。

\subsection{角动量守恒}

腾空阶段仅重力作用，角动量守恒。与二体模型一致，此处采用起跳时总角动量为零的近似假设（物理依据见第~2.3~节）：
\begin{equation}
L_G \equiv I_T\omega_T + I_A\omega_A + I_L\omega_L
       + \sum_{i\in\{T,A,L\}} m_i\bigl[(\mathbf{r}_i-\mathbf{r}_G)\times(\mathbf{v}_i-\mathbf{v}_G)\bigr]_z = 0. \label{eq:LG_three}
\end{equation}

定义相对位置 $\boldsymbol{\rho}_i \equiv \mathbf{r}_i - \mathbf{r}_G$（$i=T,A,L$）。由式\eqref{eq:rT}--\eqref{eq:rG_three}直接计算：
\begin{align}
\boldsymbol{\rho}_T &= (d_T - \beta_T)\hat{\mathbf{u}}_T - \beta_A\hat{\mathbf{u}}_A - \beta_L\hat{\mathbf{u}}_L, \label{eq:rhoT} \\
\boldsymbol{\rho}_A &= (L_T - \beta_T)\hat{\mathbf{u}}_T + (d_A - \beta_A)\hat{\mathbf{u}}_A - \beta_L\hat{\mathbf{u}}_L, \label{eq:rhoA} \\
\boldsymbol{\rho}_L &= -\beta_T\hat{\mathbf{u}}_T - \beta_A\hat{\mathbf{u}}_A + (d_L - \beta_L)\hat{\mathbf{u}}_L. \label{eq:rhoL}
\end{align}

同理计算相对速度 $\boldsymbol{\nu}_i \equiv \mathbf{v}_i - \mathbf{v}_G$：
\begin{align}
\boldsymbol{\nu}_T &= (d_T - \beta_T)\omega_T(\hat{\mathbf{k}}\times\hat{\mathbf{u}}_T)
                   - \beta_A\omega_A(\hat{\mathbf{k}}\times\hat{\mathbf{u}}_A)
                   - \beta_L\omega_L(\hat{\mathbf{k}}\times\hat{\mathbf{u}}_L), \label{eq:nuT} \\
\boldsymbol{\nu}_A &= (L_T - \beta_T)\omega_T(\hat{\mathbf{k}}\times\hat{\mathbf{u}}_T)
                   + (d_A - \beta_A)\omega_A(\hat{\mathbf{k}}\times\hat{\mathbf{u}}_A)
                   - \beta_L\omega_L(\hat{\mathbf{k}}\times\hat{\mathbf{u}}_L), \label{eq:nuA} \\
\boldsymbol{\nu}_L &= -\beta_T\omega_T(\hat{\mathbf{k}}\times\hat{\mathbf{u}}_T)
                   - \beta_A\omega_A(\hat{\mathbf{k}}\times\hat{\mathbf{u}}_A)
                   + (d_L - \beta_L)\omega_L(\hat{\mathbf{k}}\times\hat{\mathbf{u}}_L). \label{eq:nuL}
\end{align}

计算叉积 $\boldsymbol{\rho}_i \times \boldsymbol{\nu}_i$ 的 $z$ 分量，利用恒等式
\[
(\hat{\mathbf{u}}_p \times (\hat{\mathbf{k}}\times\hat{\mathbf{u}}_q))_z = \hat{\mathbf{u}}_p\cdot\hat{\mathbf{u}}_q = \cos(\theta_p - \theta_q).
\]

记 $c_{TA} \equiv \cos(\theta_T-\theta_A)$、$c_{TL} \equiv \cos(\theta_T-\theta_L)$、$c_{AL} \equiv \cos(\theta_A-\theta_L)$。对每个刚体展开 $\boldsymbol{\rho}_i \times \boldsymbol{\nu}_i$，得到 $3\times3$ 项的线性表达式。按 $\omega_T$、$\omega_A$、$\omega_L$ 合并同类项，最终将式\eqref{eq:LG_three}化为：
\begin{equation}
\boxed{A_T(\boldsymbol{\theta})\,\omega_T \;+\; A_A(\boldsymbol{\theta})\,\omega_A \;+\; A_L(\boldsymbol{\theta})\,\omega_L \;=\; 0}, \label{eq:three_body_am}
\end{equation}
其中系数函数 $A_i(\boldsymbol{\theta})$ 仅依赖当前姿态角度：
\begin{align}
A_T &= I_T + K_{TT} + K_{TT}^{(TA)} c_{TA} + K_{TT}^{(TL)} c_{TL}, \label{eq:AT} \\
A_A &= I_A + K_{AA} + K_{AA}^{(TA)} c_{TA} + K_{AA}^{(AL)} c_{AL}, \label{eq:AA} \\
A_L &= I_L + K_{LL} + K_{LL}^{(TL)} c_{TL} + K_{LL}^{(AL)} c_{AL}. \label{eq:AL}
\end{align}

这里 $K_{ii}$ 为常数项（与角度无关），$K_{ii}^{(pq)}$ 为 $\cos(\theta_p-\theta_q)$ 的系数。两者均仅为质量与几何参数的函数，完整显式表达式见附录~\ref{app:coeffs}。

\subsection{欠定约束与角速度关系}

方程\eqref{eq:three_body_am} 是一个关于三个角速度的标量约束。相比于二体模型，三体模型保留了两个独立的关节控制自由度，系统是欠定的。需要继续添加手臂角速度和腿部角速度关于时间的方程。

将式\eqref{eq:three_body_am}改写为：
\begin{equation}
\boxed{\omega_T = -\frac{A_A}{A_T}\,\omega_A \;-\; \frac{A_L}{A_T}\,\omega_L}. \label{eq:mechanical_connection}
\end{equation}

它表明：躯干的旋转角速度由手臂和腿的角速度线性给出。

\subsection{以相对角速度表示}

为便于与二体模型对比，引入相对于躯干的相对角度：
\begin{align}
\varphi_{TA} &\equiv \theta_A - \theta_T, \quad \dot\varphi_{TA} = \omega_A - \omega_T, \label{eq:phiTA} \\
\varphi_{LT} &\equiv \theta_L - \theta_T, \quad \dot\varphi_{LT} = \omega_L - \omega_T. \label{eq:phiLT}
\end{align}
其中 $\dot\varphi_{LT}$ 即为前文二体模型中的"收腿角速度" $\dot\phi$，$\dot\varphi_{TA}$ 为"摆臂角速度"。

将 $\omega_A = \omega_T + \dot\varphi_{TA}$、$\omega_L = \omega_T + \dot\varphi_{LT}$ 代入式\eqref{eq:three_body_am}，整理得：
\begin{equation}
\boxed{\omega_T = -\frac{A_A\,\dot\varphi_{TA} \;+\; A_L\,\dot\varphi_{LT}}{A_T + A_A + A_L}}. \label{eq:omega_T_relative}
\end{equation}

该式是二体模型中式\eqref{eq:dt1}的三体推广。形式上的对应如下表：

\begin{table}[h]
\centering
\caption{二体与三体模型 $\omega_T$ 表达式的结构对比}
\begin{tabular}{lcc}
\toprule
& 二体模型 & 三体模型 \\
\midrule
分子 & $(I_2 + \mu d_2^2 - \mu d_1 d_2\cos\phi)\,\dot\phi$ & $A_A\dot\varphi_{TA} + A_L\dot\varphi_{LT}$ \\
分母 & $I_1+I_2+\mu(d_1^2+d_2^2)-2\mu d_1 d_2\cos\phi$ & $A_T+A_A+A_L$ \\
控制输入 & 仅 $\dot\phi$（收腿） & $\dot\varphi_{TA}$（摆臂）+ $\dot\varphi_{LT}$（收腿） \\
\bottomrule
\end{tabular}
\end{table}

当手臂相对于躯干固定（$\dot\varphi_{TA} = 0$），即JRA情形，式\eqref{eq:omega_T_relative}简化为
\[
\omega_T = -\frac{A_L}{A_T + A_A + A_L}\,\dot\varphi_{LT},
\]
此时 $A_T + A_A$ 对应于当前姿态下"躯干+固定双臂"共同产生的角动量系数，系数形式与二体结果同构。进一步令 $m_A \to 0$、$I_A \to 0$（即消除手臂自由度），并将剩余上身参数重新识别为二体模型中的上身参数，可退化为二体模型的式\eqref{eq:dt1}（验证细节见附录~\ref{app:coeffs}）。

\subsection{落地前伸量}

三刚体模型中，脚的位置（刚体 $L$ 末端）：
\begin{equation}
\mathbf{r}_B = \mathbf{r}_O + l_L \hat{\mathbf{u}}_L,
\qquad x_B = x_O + l_L \cos\theta_L.
\end{equation}

由式\eqref{eq:rG_three}，总质心水平坐标：
\begin{equation}
x_G = x_O + \frac{m_T d_T \cos\theta_T + m_A (L_T \cos\theta_T + d_A \cos\theta_A) + m_L d_L \cos\theta_L}{M}.
\end{equation}

因此落地前伸量（脚超前质心的水平距离）：
\begin{equation}
\boxed{\Delta x = l_L\cos\theta_L - \frac{m_T d_T \cos\theta_T + m_A (L_T \cos\theta_T + d_A \cos\theta_A) + m_L d_L \cos\theta_L}{M}}. \label{eq:Dx_three_body}
\end{equation}

\subsection{数值求解}

三体模型的数值求解包含两个"控制输入"：
\begin{enumerate}
\item 摆臂规律：给定手臂相对于躯干的角速度 $\dot\varphi_{TA}(t)$。
\item 收腿规律：给定腿相对于躯干的角速度 $\dot\varphi_{LT}(t)$。
\end{enumerate}

求解流程与二体模型欧拉法类似：
\begin{enumerate}
\item 由式\eqref{eq:omega_T_relative} 代入当前姿态计算 $\omega_T(t)$。
\item 积分更新 $\theta_T$、$\varphi_{TA}$、$\varphi_{LT}$，进而得到 $\theta_A$、$\theta_L$。
\item 代入式\eqref{eq:Dx_three_body} 计算 $\Delta x$。
\end{enumerate}

对应的 Python 代码仅是二体代码的直接扩展：在每一步由 $A_T$、$A_A$、$A_L$（见式\eqref{eq:AT}--\eqref{eq:AL}，系数由附录~\ref{app:coeffs} 提供）计算 $\omega_T$，再欧拉步进。

\medskip
\noindent\textbf{算例演示。} 为展示框架的具体计算能力，取如下控制输入：收腿角速度维持 $\dot\varphi_{LT}=180^\circ/\text{s}$（与二体模型一致）；摆臂采用分段恒速策略——前 $0.25$\,s 以 $\dot\varphi_{TA}=+180^\circ/\text{s}$ 前摆（延续起跳阶段的前摆动量），后 $0.25$\,s 以 $\dot\varphi_{TA}=-180^\circ/\text{s}$ 后摆（为落地姿态做准备）。初始姿态 $\theta_T(0)=60^\circ$，$\varphi_{TA}(0)=30^\circ$，$\varphi_{LT}(0)=-155^\circ$。数值积分（步长 $0.001$\,s）得到表~\ref{tab:three_body_results} 的对比结果。

\begin{table}[h]
\centering
\caption{三刚体模型典型算例结果对比}
\label{tab:three_body_results}
\begin{tabular}{lccc}
\toprule
& 二体模型 & 三体（固定臂） & 三体（摆臂） \\
\midrule
落地 $\theta_T$ ($^\circ$) & $\approx 14$ & 23.0 & 22.9 \\
落地 $\theta_L$ ($^\circ$) & $\approx -52$ & $-42.0$ & $-42.1$ \\
落地 $\Delta x$ (m) & 0.298 & 0.383 & 0.382 \\
\bottomrule
\end{tabular}
\end{table}

三体固定臂情形（$\dot\varphi_{TA}=0$）的 $\Delta x=0.383$\,m 高于二体模型的 $0.298$\,m，这是因为手臂从躯干刚体中分离为独立刚体后改变了系统的总质心位置和角动量分配，即使手臂固定不动，其对 $\Delta x$ 的贡献也与并入躯干的杆模型假设不同。摆臂情形（先向前、后向后）的 $\Delta x=0.382$\,m 与固定臂结果几乎相同（差距 $<0.1$\,cm），并非计算误差——它揭示了腾空阶段角动量守恒的根本约束：在纯惯性空间中，手臂前摆所"贡献"的角动量必然伴随躯干与腿部的反向转动，后半程的后摆又将这一效果基本抵消，因此手臂在腾空阶段单独无法显著改变落地前伸量。

这一结论与 Ashby \& Delp (2006) 的仿真发现一致：手臂对跳跃成绩的贡献中约 80\% 来自起跳阶段（地面反力提供外力矩，手臂摆动可有效改变系统总能量和起跳姿态），仅有少量来自腾空姿态的微调（$\Delta x$ 增加约 $2$\,cm）。三刚体模型从第一性原理定量展示了：在腾空角动量守恒这一硬约束下，无论手臂轨迹如何设计，单靠腾空阶段的手臂运动均无法对落地前伸量产生大幅提升。由此自然引出下一阶段的研究方向——将起跳阶段（含地面反力）与腾空阶段统一建模，将 $\dot\varphi_{TA}(t)$ 和 $\dot\varphi_{LT}(t)$ 作为控制变量，构成一个可数值求解的最优控制问题。本文给出的 $A_T$、$A_A$、$A_L$ 精确系数表达式为该方向提供了完整的力学基础。

\section{结论}

本文从简到繁、逐步深入地建立了立定跳远腾空阶段的力学模型。首先，二刚体模型利用角动量守恒约束，定量描述了腾空阶段收腿对落地前伸量的贡献。模型预测的落地前伸量 $\Delta x = 0.298$\,m 与实验限制摆臂数据（$0.29$\,m）偏差 $\ll 3\%$，表明二刚体模型准确地捕捉了腾空阶段姿态动力学的主要特征。其次，根据 Ashby \& Heegaard (2002) 的实验观测，手臂摆动对跳跃成绩的提升主要体现在起跳阶段（起跳速度增加约 13\%，贡献约 80\% 的距离增益），腾空姿态的改善（落地前伸量增加约 2 cm）幅度较小。最后，三刚体模型的推广从解析建模给出了手臂与躯干、下肢之间角动量分配的定量表达式，并从第一性原理展示了腾空角动量守恒对手臂摆动效果的根本限制，为后续通过最优控制方法研究摆臂与收腿的协同策略提供了动力学基础。

\appendix
\section{二刚体模型参数的估算依据}

正文表~\ref{tab:two_body_params} 中的转动惯量和收腿角速度按以下方式估算。

\textbf{转动惯量。}两刚体分别近似为均匀杆，绕自身质心的转动惯量为 $I = \frac{1}{12}mL^2$。
\begin{itemize}
  \item \textbf{刚体 1（躯干+头+臂）}：质心到髋关节距离 $d_1=0.35$\,m，均匀杆质心在中点，杆长 $L_1 = 2d_1 = 0.70$\,m。则
    \[
    I_1 = \frac{1}{12}m_1 L_1^{\,2}
        = \frac{1}{12}\times 44 \times 0.70^2
        = 1.80\;\text{kg\,m$^2$}
        \approx 1.8\;\text{kg\,m$^2$}.
    \]
  \item \textbf{刚体 2（双腿+足）}：腿总长 $l_2 = 0.95$\,m。则
    \[
    I_2 = \frac{1}{12}m_2 l_2^{\,2}
        = \frac{1}{12}\times 28 \times 0.95^2
        = 2.11\;\text{kg\,m$^2$}
        \approx 2.1\;\text{kg\,m$^2$}.
    \]
\end{itemize}

\textbf{收腿角速度。}起跳时 $\phi(0) = -155^\circ$，落地时腿需收至躯干前方。取恒速 $\dot\phi = 180^\circ/\mathrm{s}$（$\approx 3.14\ \mathrm{rad/s}$），总幅度 $\Delta\phi = 90^\circ$，落地 $\phi(T) = -65^\circ$，足以使脚到达质心前方（$\Delta x > 0$）。该值与腾空阶段髋关节屈曲速率的典型观测一致（$150$--$250^\circ/\mathrm{s}$）。实际收腿非恒速，此处为平均近似。

\section{三刚体角动量系数的完整展开}
\label{app:coeffs}

本节给出正文式\eqref{eq:AT}--\eqref{eq:AL}中所有系数 $K_{ij}^{(\,\cdot\,)}$ 的显式形式。这些表达式完全由$\beta_T,\beta_A,\beta_L$和原始质量、几何参数确定，是 $\boldsymbol{\rho}_i \times \boldsymbol{\nu}_i$ 各项系统展开并合并同类项的结果。

首先回顾约化距离的定义（式\eqref{eq:beta_def}）：
\begin{equation}
\beta_T = \frac{m_T d_T + m_A L_T}{M},\quad
\beta_A = \frac{m_A d_A}{M},\quad
\beta_L = \frac{m_L d_L}{M},\quad
M = m_T + m_A + m_L.
\end{equation}

\subsection*{A. 常数项（不含角度余弦）}

\begin{align}
K_{TT} &= m_T(d_T - \beta_T)^2 + m_A(L_T - \beta_T)^2 + m_L\beta_T^2, \label{eq:KTT} \\
K_{AA} &= (m_T + m_L)\beta_A^2 + m_A(d_A - \beta_A)^2, \label{eq:KAA_const} \\
K_{LL} &= (m_T + m_A)\beta_L^2 + m_L(d_L - \beta_L)^2. \label{eq:KLL_const}
\end{align}

\subsection*{B. $\cos(\theta_T-\theta_A)$ 的系数}

\begin{align}
K_{TT}^{(TA)} = K_{AA}^{(TA)} &= -\,m_T\beta_A(d_T-\beta_T)
                                 + m_A(d_A-\beta_A)(L_T-\beta_T)
                                 + m_L\beta_A\beta_T. \label{eq:K_TA}
\end{align}

\subsection*{C. $\cos(\theta_T-\theta_L)$ 的系数}

\begin{align}
K_{TT}^{(TL)} = K_{LL}^{(TL)} &= -\,m_T\beta_L(d_T-\beta_T)
                                 - m_A\beta_L(L_T-\beta_T)
                                 - m_L\beta_T(d_L-\beta_L). \label{eq:K_TL}
\end{align}

\subsection*{D. $\cos(\theta_A-\theta_L)$ 的系数}

\begin{align}
K_{AA}^{(AL)} = K_{LL}^{(AL)} &= m_T\beta_A\beta_L
                                 - m_A\beta_L(d_A-\beta_A)
                                 - m_L\beta_A(d_L-\beta_L). \label{eq:K_AL}
\end{align}

\subsection*{E. 系数对称性说明}

注意 $K_{TT}^{(TA)} = K_{AA}^{(TA)}$、$K_{TT}^{(TL)} = K_{LL}^{(TL)}$、$K_{AA}^{(AL)} = K_{LL}^{(AL)}$。该对称性来源于角动量表达式中各刚体通过总质心相互耦合的二次型结构——各刚体通过总质心相互耦合，使得 $\omega_i$ 与 $\omega_j$ 的交叉系数相等。利用这些对称性可将正文式\eqref{eq:AT}--\eqref{eq:AL}中的 $9$ 个系数降至仅 $6$ 个独立参数。

\subsection*{F. 退化验证}

令 $m_A \to 0$、$I_A \to 0$（消除手臂自由度），则 $\beta_A \to 0$，且 $\beta_T \to m_T d_T / (m_T+m_L)$。此时：
\begin{itemize}
\item $K_{AA}$ 及 $K_{AA}^{(TA)}$、$K_{AA}^{(AL)}$ 均趋于零，故 $A_A \to 0$。
\item $K_{TT}^{(TA)}$、$K_{AA}^{(TA)}$ 趋于零；$K_{AA}^{(AL)}$、$K_{LL}^{(AL)}$ 也趋于零。
\end{itemize}

式\eqref{eq:three_body_am} 退化为：
\[
\bigl[I_T + K_{TT} + K_{TT}^{(TL)}c_{TL}\bigr]\omega_T
+ \bigl[I_L + K_{LL} + K_{LL}^{(TL)}c_{TL}\bigr]\omega_L = 0.
\]

若进一步将剩余上身刚体参数识别为二体模型的上身参数，即取对应关系 $m_T = m_1$、$I_T=I_1$、$m_L = m_2$、$I_L=I_2$、$d_T = d_1$、$d_L = d_2$，并注意 $c_{TL} = \cos(\theta_T-\theta_L) = \cos(\theta_1-\theta_2)$，则将 $K_{TT}$、$K_{LL}$、$K_{TT}^{(TL)}$ 代入展开，所得系数与正文二体模型式\eqref{eq:angular_momentum_exact}中的对应项系数完全一致，验证了三体模型的自洽性。

\end{document}
